%%%%%%%%%%  scan idx 613  |  folio 600  %%%%%%%%%%
\[
\begin{tikzcd}
Y & \coprod_j Y_j \arrow[l]\\
X \arrow[u] & \coprod_j \Sigma_j \arrow[l] \arrow[u] & ,
\end{tikzcd}
\]
\[
\begin{tikzcd}
\underline{\Pic}_{Y/S} \arrow[r,"\sgacircled{1}"] \arrow[d,"\sgacircled{2}"'] &
\underline{\Pic}_{\coprod_jY_j/S}=\prod_j\underline{\Pic}_{Y_j/S} \arrow[d,"\sgacircled{4}"]\\
\underline{\Pic}_{X/S} \arrow[r,"\sgacircled{3}"] &
\underline{\Pic}_{\coprod_j\Sigma_j/S}=\prod_j\underline{\Pic}_{\Sigma_j/S}.
\end{tikzcd}
\]
Et pour établir II, il suffit de voir que, quitte à restreindre \(S\), \(\sgacircled{1}\), \(\sgacircled{3}\) et \(\sgacircled{4}\) sont quasi-affines.

Considérons \(\sgacircled{4}\) d'abord. En remplaçant \(Y\) par \(Y_j\) et \(X\) par \(\Sigma_j\), on se ramène au cas \(Y\) intègre sur \(S\). Comme \(f\) est surjectif, il existe l'un des \(X_i\), soit \(X_0\), qui domine \(Y\), et on a les deux suites
\[
  X_0\hookrightarrow \coprod X_i\longrightarrow Y,
\]
\[
  \Pic_{X_0/S}\longleftarrow \Pic_{\coprod X_i/S}\longleftarrow \Pic_{Y/S}.
\]
Si l'on sait donc que les \(\underline{\Pic}_{X_i/S}\) donc aussi \(\Pic_{\coprod X_i/S}\) sont représentables, et le morphisme composé est quasi-affine, on conclut que \(\sgacircled{4}\) est quasi-affine. Compte tenu de la démonstration de 1.2 on est donc ramené pour \(\sgacircled{4}\) à démontrer II'. \uline{Le théorème} 1.1 \uline{lorsque} \(X\) \uline{et} \(Y\) \uline{sont intègres sur} \(S\).

Pour \(\sgacircled{1}\) et \(\sgacircled{3}\), on doit traiter le cas d'un morphisme \(\coprod X_i\to X\) où \(X\) est réduit, les \(X_i\) sont les composantes irréductibles de \(X\) et sont intègres sur \(S\). Or on peut supposer \(X\) à fibres connexes. En effet, soit \(\eta\) le point géométrique de \(S\). Alors \(X_\eta\) a comme composantes irréductibles les \((X_i)_\eta\) et, quitte à restreindre \(S\), on peut supposer que \(X_i\cap X_j=\varnothing\) si et seulement si
